\section{Modelling the traffic}
\label{sec:traffic}

Much of what a fleet simulator tells you depends on what arrives at it. The
textbook default is Poisson arrivals: jobs appearing independently of one another
at a constant average rate, with exponentially distributed gaps. Poisson traffic
is smooth in a specific sense --- average it over a longer window and its
relative variability shrinks predictably. Measured computer traffic does not do
that. It is bursty at every timescale you look at, and averaging does not smooth
it out \cite{paxsonfloyd1995}. A study built on Poisson arrivals understates
queue variance, and understates it most exactly where the answer matters, near
saturation. fleetsim therefore runs one arrival process per class and superposes
them (Table~\ref{tab:arrivals}).

\begin{table}[t]
\centering\small
\caption{Arrival process per job class, with shipped defaults. Pretraining is
the one surviving Poisson case: rare planned launches really are independent,
human-initiated events.}
\label{tab:arrivals}
\begin{tabular}{lll}
\toprule
Class & Process & Defaults \\
\midrule
Pretrain    & homogeneous Poisson              & \texttt{rate\_per\_week: 2} \\
Finetune    & MMPP-2 $\times$ diurnal envelope & \texttt{rate\_ratio: 4}, \texttt{burst\_frac: 0.25}, \texttt{switch\_tau: 2d} \\
Eval        & Hawkes on an NHPP baseline       & \texttt{branching: 0.4}, \texttt{kernel\_tau: 15m} \\
Inference   & service replica curve            & --- \\
Best-effort & closed-loop backlog              & \texttt{concurrency: 4}, \texttt{think\_time: fixed[0s]} \\
Retry storm & near-critical Hawkes (preset)    & \texttt{branching: 0.9}, \texttt{kernel\_tau: 2m} \\
\bottomrule
\end{tabular}
\end{table}

The shared day-and-week machinery is a non-homogeneous Poisson process whose log
rate is a sum of harmonics at 24\,h and 168\,h. The exponential keeps the rate
positive, and the constant term is derived rather than configured, so changing
the shape of a day does not change the amount of work. The shipped daily preset
is a two-harmonic least-squares fit to \texttt{fleetsim}'s own v0.1 step curve,
which was itself built from the Helios description of a night trough with dips
at 12:00 and 18:00 \cite{helios2021}. The provenance matters because the fit
loses something: two harmonics smooth those two dips away, leaving a trough near
$0.59\times$ at 03:00 and a single peak near $1.27\times$ at 10:15. The shipped
preset is therefore continuous with the previous release's behaviour rather than
a fit to a published curve, and it no longer reproduces the midday and evening
dips Helios reports. It is sampled by Lewis--Shedler thinning --- propose
arrivals at a constant upper-bound rate, then keep each with probability equal
to the true rate over that bound.

Finetuning sits on a two-state Markov-modulated Poisson process: a hidden switch
between a quiet and a busy regime, reproducing crunch weeks that a single rate
curve cannot. It is configured in observable terms --- mean rate, burst fraction,
switch correlation time, rate ratio --- and solved to state rates in closed form.
Those rates must be scaled to the \emph{peak} of the diurnal envelope before
thinning; left unscaled they silently generate about 60\,\% of the configured
load.

Evaluation traffic uses a Hawkes process. Self-exciting means one submission
makes further submissions more likely for a while afterwards: a failed eval is
re-run, a user iterates. Configuration exposes the branching ratio --- mean
children per arrival, stationary below 1 --- and the kernel decay time; sampling
uses Ogata thinning, sequential but $O(1)$ per event because the excitation term
decays exponentially.

Sizes are per-class distributions over powers of two. Evals are 55\,\%
single-chip: inside the 50--90\,\% single-GPU range Helios reports, and
deliberately below Philly's finding that over 82\,\% of jobs are single-GPU
\cite{philly2019}. TPU clusters use the published histogram of \emph{slice}
sizes \cite{tpuv4isca2023} --- a TPU slice being a contiguous geometric block of
chips carved out of a pod, a different thing from the quota slices of
\S\ref{sec:helios}. Uniform-over-exponents sizing remains available but is not
the default, since it overweights giants by 5--10$\times$ and fabricates
queueing.

Job durations are drawn from a lognormal shape, configured by naming its
quantiles (median and 90th percentile) rather than its abstract parameters.
Pretraining additionally gets a \emph{heavy tail}: a small probability of a run
lasting far longer than the median, decaying as a power law rather than dying
off exponentially. How heavy a tail is, is summarised by its \emph{tail index};
a smaller index means a heavier tail. Pretraining gets the only heavy tail
because resource-hours are size $\times$ duration, and a product of two
quantities is as heavy-tailed as the heavier of them --- it inherits the
\emph{smaller} of the two indices. Attaching the tail to the one class with
genuinely open-ended runs is therefore intended to suffice for the Borg-like
resource-hour tail \cite{borgtng2020}, without fabricating implausibly long
evaluation jobs. Intended, not confirmed: the distributional self-check that
would measure the generated chip-hour tail is a documented target in the traffic
specification, not a shipped test.

Tenants come from a \emph{finite} Zipf over ranks. An exponent of 1.19
reproduces Alibaba PAI's top 5\,\% of users submitting about 77\,\% of
instances \cite{alibabapai2022}, at a population of 1{,}300 tenants; the shipped
default is 1.2. The class mix follows Acme \cite{acme2024}.

These parameters are sourced but not all fitted, and the repository says so: the
Markov-modulated and Hawkes defaults are order-of-magnitude priors borrowed from
grid and LLM-serving analogies, not constants fitted to an ML fleet. There is no
explicit long-range-dependence generator, tenants are marks on jobs rather than
agents, and within a class size and duration are drawn independently.

\section{Validation}
\label{sec:validation}

\subsection{What ``validated'' means here, and the two kinds of check}
\label{sec:val-kinds}

A simulator can be internally consistent and still be wrong about the world.
Validation here means comparing fleetsim's output against numbers a published
study measured on a real cluster, and two kinds of comparison are possible. A
\textbf{policy-effect check} runs two scheduling policies over the same trace and
compares the \emph{ratio} between them against the ratio the paper reported. A
\textbf{distribution-match check} asks whether a converted trace reproduces a
published shape, such as the fraction of jobs that ended successfully; that
mostly tests the converter rather than the scheduler, and in the Philly case
involves no simulation at all.

The ratio is the more robust test, and the reason deserves two extra sentences.
Suppose the fleet we reconstructed is 10\,\% smaller than the real one. Both
policies then run on a fleet 10\,\% too small, both produce waits that are too
long, and much of that shared error divides out when one is taken over the other.
An absolute number has nothing to divide it out; it carries every reconstruction
error straight through.

Replays never match exactly, for structural reasons. Take our own replay first.
Helios publishes its analysis window --- 2020-09-01 to 2020-09-26 --- and we
replay exactly that, so the window is not a source of error here. Two others
remain. Timestamps carry no timezone, so a daily cycle cannot be aligned with
confidence. And the original scheduler made placement decisions nobody recorded,
which is the residual \S\ref{sec:replay-found} spends most of its length on.

Where a released trace and its paper's analysis window differ, the problem is
worse. Philly released 117{,}325 jobs over 137 days while the paper analysed
96{,}260 over roughly 75, and that window is not precisely published, so no
reconstruction of it can be checked. The tolerance bands below are sized to
absorb all of those residuals; they are not calibrated uncertainty intervals.

\subsection{Replaying a real trace: the Helios experiment}
\label{sec:helios}

The flagship replay targets Hu et al.'s Helios study \cite{helios2021}, which
characterizes four production GPU clusters --- Venus, Earth, Saturn and Uranus
--- over 2020-09-01 to 2020-09-26, GPU jobs only. The trace is a
36{,}437{,}672-byte archive (343{,}069{,}364 bytes uncompressed) under CC-BY-4.0;
fleetsim fetches it, verifies it by size, and replays all four clusters in about
four minutes.

Each cluster is partitioned into \emph{virtual clusters}: quota slices, each a
group's private share of the machines, scheduled independently. The reference
simulator runs one scheduler per slice, whereas fleetsim schedules a single
global pool, so the replay is a harness rather than an engine change --- one
simulation per slice, on a fleet sized to that slice's node count at 8 GPUs per
node, aggregated job-weighted to cluster totals. The converter writes each job's
recorded duration into both its duration and its runtime estimate, making
shortest-job-first an \emph{oracle}: a perfect estimate, and so a strict upper
bound on what a duration-\emph{predicting} policy could achieve.

Every claim below holds only under three modelling choices, all options a reader
could disagree with: a strict blocking scan, September-max per-slice sizing, and
\texttt{consolidate} placement. Two of them were not chosen but forced ---
Table~\ref{tab:corrections} records the assumptions the real data disproved. A
strict scan means the queue head blocks everyone behind it until it can run,
which is what head-of-line blocking means; both policies share the scan mode, so
shortest-job-first's advantage remains purely its ordering. September-max sizing
drops zero jobs, and the test asserts that, but it is not unambiguously more
faithful and the repository says so: peak quota available from day one
over-provisions early in the month and under-counts early-month queueing. The
converter's own default remains the snapshot.

\begin{table}[t]
\centering\small
\caption{Two modelling assumptions the validation plan made before the data was
touched, and what the real trace showed.}
\label{tab:corrections}
\begin{tabular}{p{0.14\textwidth}p{0.30\textwidth}p{0.46\textwidth}}
\toprule
Choice & The plan assumed & What the trace showed \\
\midrule
Scan mode &
Non-blocking: a job the scheduler cannot place is skipped and later jobs
proceed. &
The reference blocks. Non-blocking collapses FIFO's queueing --- ratios near
$1.4\times$ with the wrong cross-cluster ordering. A strict scan flips the
result from non-reproducing to reproducing. \\
Per-slice capacity &
A fixed quota snapshot taken on 1 September. &
Quota drifts mid-month: one Uranus slice grows 208 $\to$ 328 GPUs, another
spins up 0 $\to$ 416, and 3{,}117 jobs sit in slices with no 1 September quota
at all, which a snapshot silently drops. It also over-congests Uranus (ratio
$3.41\times$), breaking the rank and the queuing-share ordering. \\
\bottomrule
\end{tabular}
\end{table}

Table~\ref{tab:helios-ratios} gives the result: all eight ratios in band. The
quantity being compared is the job completion time, or JCT --- the wall-clock
time from when a job was submitted to when it finished, so queue wait plus run
time. The test also asserts that no job was dropped, that every job reached a
terminal state, that Saturn has the highest and Uranus the lowest JCT ratio, and
that the queuing share of JCT orders Saturn $>$ Venus $>$ Earth $>$ Uranus in
fleetsim (90.1, 79.1, 73.0, 42.2\,\%) as in the published table (89.7, 81.8,
69.3, 42.5\,\%). It runs in roughly four minutes on a laptop against the cached
trace: the repository's own timings are 250\,s, ``about 260\,s'' and 266\,s in
three different files, which is the accuracy a wall-clock number of this kind
deserves. The bands are wide on purpose, and were deliberately \emph{not}
tightened even though every JCT ratio now lands within 5\,\% of published, for
the two reasons in \S\ref{sec:val-kinds}. Nor is the queuing column an exact
reproduction: Earth measures $10.67\times$ against a published $16.4\times$,
about 35\,\% low, and passes only because the band runs to $25\times$. The band
does have teeth --- SJF replayed against itself gives 1.0 and fails it.

\begin{table}[t]
\centering\small
\caption{Helios replay, four clusters, real trace, shipped configuration (strict
scan, September-max sizing, \texttt{consolidate} placement). The
\texttt{fleetsim} column is the claim. The \texttt{first\_fit} column is the same
replay under the engine's \emph{default} placer, shown so the size of the
placement effect is visible rather than asserted. The \emph{In band} columns
judge the \texttt{fleetsim} column only. Bands are $[1.3, 8]\times$ on the
completion-time ratio and $[3, 25]\times$ on the queuing ratio; $\dagger$ marks
the one \texttt{first\_fit} value that falls outside its band, the deviation
\S\ref{sec:replay-found} diagnoses. The mean absolute errors quoted in the text
(3.4\,\% for \texttt{fleetsim}, 27.4\,\% for \texttt{first\_fit}) are computed on
full-precision ratios; the rounded values shown here give 3.2\,\% and 27.5\,\%.}
\label{tab:helios-ratios}
\begin{tabular}{lrrcrrrcr}
\toprule
& \multicolumn{4}{c}{FIFO/SJF mean job-completion-time ratio} & \multicolumn{4}{c}{FIFO/SJF mean-queuing ratio} \\
\cmidrule(lr){2-5}\cmidrule(lr){6-9}
Cluster & Published & fleetsim & In band & \texttt{first\_fit} & Published & fleetsim & In band & \texttt{first\_fit} \\
\midrule
Saturn & 6.59 & 6.87 & yes & 8.75$^\dagger$ & 18.5 & 19.55 & yes & 22.91 \\
Venus  & 3.07 & 3.21 & yes & 4.21           & 5.68 & 7.78  & yes & 9.97  \\
Earth  & 2.87 & 2.95 & yes & 2.11           & 16.4 & 10.67 & yes & 5.73  \\
Uranus & 1.49 & 1.51 & yes & 1.69           & 4.51 & 5.08  & yes & 6.10  \\
\bottomrule
\end{tabular}
\end{table}

\subsection{What the replay found}
\label{sec:replay-found}

\paragraph{The explanation we had was impossible.}
The previous version of the suite had one out-of-band number: Saturn's JCT ratio
at $8.75\times$ against a published $6.59\times$. It also had an explanation ---
that first-fit placement fragments Saturn's large multi-node jobs more than the
reference implementation's consolidating placer. That explanation was wrong, and
provably rather than merely unsupported. The replay fleet is single-level: every
node is a direct child of the cluster root, replay jobs carry no locality
constraint and no segment structure, and no crossing penalty is configured, so
\emph{which} free nodes a multi-node job takes cannot change any outcome at all.
The sizes agree: Saturn's largest job is 200 GPUs, and the quota slice carrying
84\,\% of the FIFO-minus-SJF gap tops out at 56 GPUs, with 0.7\,\% of its jobs
larger than one node. It was never a large-job effect.

\paragraph{The real mechanism: sub-node stranding.}
One 1-GPU job landing on an empty 8-GPU node is enough. That node now has seven
free GPUs and an owner, and the allocator treats any node with an owner as
ineligible for a whole-node request --- so a 16-GPU gang cannot touch it. Do
that on enough nodes and the fleet reports plenty of free chips while no gang
can start.

That is what first-fit does here. Saturn is 70.8\,\% single-GPU jobs, 97.1\,\%
in its dominant quota slice. First-fit scans nodes by ascending identifier with
no preference for nodes already partly used, so it opens a fully free node for a
1-GPU job whenever the lower-numbered ones are full, and free capacity
accumulates as 1--7-GPU remainders no job of 8 GPUs or more can claim. The
whole-node rule itself is an approximation whose error was quantified rather
than assumed: Helios's own node count equals $\lceil \mathit{gpus}/8 \rceil$ for
99.66\,\% of windowed multi-GPU jobs, and every exception spreads a job over
\emph{more} nodes than the ceiling, never fewer --- ``never fewer'' being the
half the rule depends on.

Instrumented runs confirmed the mechanism: 100\,\% of measured blocked-idle
chip-seconds come from queue heads needing whole nodes, sub-node heads
contribute 0.0\,\%, and 87.7\,\% of free chips during blocks sit on partly-used
nodes. This happens at an \emph{offered load} of 94.4\,\% --- offered load being
the arriving work as a fraction of the fleet's capacity to do it. Near 100\,\%
queues become extremely sensitive: a few percent of capacity made unusable does
not lengthen the queue by a few percent, it lengthens it by a multiple
\cite{kingman1961}. Those instrumentation figures came from throwaway code not
in the tree --- the diagnostic work that found the mechanism, not numbers the
suite regenerates.

\paragraph{What the fix did and did not buy.}
Switching the sub-node rule to tightest-fit packing --- the \texttt{consolidate}
policy, which fills partly-used nodes before opening empty ones --- closes most
of the gap. The four-cluster mean absolute error against the published ratios
falls from 27.4\,\% to 3.4\,\%, and Saturn's mean FIFO completion time moves
from 75{,}329\,s to 55{,}978\,s against a published 55{,}984\,s. That 0.01\,\%
is luck, not precision --- the tie-break probe below moves the same quantity by
17\,\%. What is supported is the band across four clusters in
Table~\ref{tab:helios-abs}, never a single cell.

\begin{table}[t]
\centering\small
\caption{Absolute quantities: reported, \emph{not} asserted --- continuous
integration does not fail if they move. \texttt{\#Queuing} counts jobs whose wait
exceeds one 60\,s scheduler round. $\ast$ marks the two quantities that moved
\emph{away} from published when the placement model was corrected.}
\label{tab:helios-abs}
\begin{tabular}{lrrrrrr}
\toprule
& \multicolumn{3}{c}{Mean FIFO job completion time (s)} & \multicolumn{3}{c}{\texttt{\#Queuing} jobs} \\
\cmidrule(lr){2-4}\cmidrule(lr){5-7}
Cluster & Published & fleetsim & $\Delta$ & Published & fleetsim & $\Delta$ \\
\midrule
Saturn & 55{,}984 & 55{,}978 & $-0.01$\,\%       & 65{,}991 & 72{,}936 & $+10.5$\,\% \\
Venus  & 64{,}702 & 55{,}714 & $-13.9$\,\%       & 15{,}336 & 13{,}624 & $-11.2$\,\%$^\ast$ \\
Earth  & 19{,}754 & 22{,}463 & $+13.7$\,\%       & 30{,}030 & 32{,}232 & $+7.3$\,\% \\
Uranus & 19{,}758 & 18{,}492 & $-6.4$\,\%$^\ast$ & 16{,}917 & 18{,}449 & $+9.1$\,\% \\
\bottomrule
\end{tabular}
\end{table}

Three qualifications belong with that improvement. One change did not fix
everything: Uranus's absolute FIFO JCT went from $+5$\,\% to $-6.4$\,\%,
overshooting past published rather than approaching it; Venus's queued-job count
drifted from $-9$\,\% to $-11.2$\,\%; and the published \emph{ranking} of
absolute JCTs, which the earlier version happened to reproduce, no longer is ---
though nobody can meaningfully reproduce it, since published Uranus
(19{,}758\,s) and Earth (19{,}754\,s) differ by 0.02\,\%, a dead tie. Only the
ranking of the \emph{ratios} is asserted. Nor are the eight absolute quantities
eight independent confirmations: mean JCT is service plus mean queueing, the
ratio is one mean JCT over another, and the queuing share is queueing over JCT,
so all are functions of the same per-cluster wait distribution and together are
close to one degree of freedom. And the new placer does not dominate the old ---
of the five Saturn quota slices carrying 97\,\% of the gap, three get
\emph{worse} under it, one by 18.5\,\%, the cluster-level win being carried by a
single slice holding 41\,\% of Saturn's jobs. The supported claim is that the policy
reproduces the reference placer's \emph{aggregate} behaviour, never that it is a
better scheduler.

\paragraph{The bands do not discriminate between placement policies.}
\texttt{spread}, a deliberately bad control policy that scatters work onto the
emptiest nodes, was measured across all four clusters and satisfies every
published-facing check above --- both bands on all four, both rank invariants,
the queuing-share ordering --- while being 45\,\% off Saturn's absolute FIFO JCT
and 99\,\% off Earth's. Passing the bands is therefore not by itself evidence
that the placement model is right, and ``in band'' must never be read as
``reproduces the paper''. The four-cluster version of that result is a
documented measurement; what the suite \emph{asserts} is the Saturn instance of
it, so a placement regression fails loudly on one cluster rather than on four.
The point of asserting it at all is that the misreading cannot then be made
silently.

What selects the placer instead, in descending weight, is: the four-cluster mean
absolute ratio error, asserted below 10\,\% and measuring 3.4\,\% against
27.4\,\% for first-fit and 29.5\,\% for the control, since aggregating four
clusters stops any one cluster's ordering noise from deciding it; a mechanism
metric --- the time-averaged count of partly-occupied healthy nodes, fleet state
rather than a queue statistic --- which orders the three policies exactly as
their FIFO completion times do; and Saturn's absolute number, which corroborates
but cannot independently discriminate. It cannot because 35.5\,\% of Saturn's
jobs share an exact submit second, and reversing the tie-break inside those
seconds moves Saturn's FIFO JCT by 17\,\% under first-fit --- enough to invert
which placer looks closer to published. That is a sensitivity probe under a
deliberately wrong arrival order, not an error bar: ascending identifier is
demonstrably the faithful one. Whether the four-cluster discriminator would
survive that perturbation is unmeasured. The probe ran on one cluster under one
scheduler --- Saturn under FIFO, in both tie-break orders --- and extending it
to all four clusters under both schedulers needs 24 more replays. An open
question, not a claim.

\subsection{Closed-form checks}
\label{sec:closed-form}

Trace replay says whether the simulator matches one real cluster; closed-form
checks say whether the engine underneath computes what queueing theory already
knows. The checks form a ladder, from pure algebra at the bottom to a real trace
at the top, and the repository calls each step a \emph{rung}; we keep its word.

Two rungs need their notation unpacked. Queueing theory labels a system by its
arrival process, its service-time distribution and its server count: M/M/c means
memoryless arrivals, memoryless service times and $c$ servers, and M/G/1 means
memoryless arrivals, a general service-time distribution and one server. For a
handful of such systems the mean wait has an exact formula --- Erlang-C for
M/M/c, Pollaczek--Khinchine for M/G/1 --- so we can build the same system in
fleetsim, run it, and check the simulated mean wait against the algebra.

Each rung in Table~\ref{tab:rungs} runs in continuous integration and is chosen
for a specific failure it would catch. Not everything planned is shipped, and
the unshipped items are close cousins of shipped ones, so they are worth naming
precisely. The M/G/1 conservation-law identity --- a statement that a certain
weighted sum of waits is the same under any non-anticipating discipline, and a
different object from the chip-hour accounting invariant in the table --- is a
documented target. So is the M/M/c generalisation of the two-tier priority rung,
whose shipped form is single-server. So are most distributional self-checks on
the traffic generator.

\begin{table}[t]
\centering\small
\caption{Closed-form and invariant rungs, and what each would catch.}
\label{tab:rungs}
\begin{tabular}{p{0.34\textwidth}p{0.56\textwidth}}
\toprule
Rung & What it would catch \\
\midrule
Erlang-C (M/M/c): mean wait and occupancy within 10\,\% of analytic &
A broken queue discipline, a mis-scaled arrival rate, or a double-counting
occupancy integral. \\
Pollaczek--Khinchine (M/G/1) under lognormal service &
Duration-sampler bias the exponential-only rung structurally cannot see, since
the formula depends on the second moment of service time (a distribution's
spread, not just its average, drives the wait). \\
Preemptive-resume priority, two tiers, \emph{one server} (M/M/1), plus a
shielding test under a saturating backlog &
Preemption-accounting errors and leakage of low-priority load into
high-priority wait. Best-effort mean wait is asserted on nothing --- it is
undefined under saturation. \\
Shortest-processing-time optimality, requiring an identical makespan &
A reordering policy that changes how much work is done rather than only when it
is done. \\
EASY backfill \cite{easybackfill1995}, carrying its documented negative case
under honest over-estimates &
An implementation that silently enforces runtime estimates, which real clusters
do not. \\
Invariants: gang atomicity, chip-hour conservation as a per-flush accounting
identity, per-flush allocation
accounting, quota caps with labelled demotions, reservation exclusivity, the
exact-microsecond crossing-penalty identity, byte-identical output for one
scenario and seed on one platform &
The bug class that yields plausible numbers from an internally inconsistent
fleet state. \\
\bottomrule
\end{tabular}
\end{table}

\subsection{What fleetsim deliberately does not reproduce}
\label{sec:anti-goals}

Table~\ref{tab:antigoals} lists the published numbers fleetsim does not try to
match. They are stated so that nobody quotes this paper for something it did not
test; none is asserted anywhere in the suite, so none can surface as a test
failure, which is precisely why they must be written down. Two rows deserve a
note: the fragmentation-delay split and fractional-GPU sharing were both targets
for the release that then spent its budget on placement, and are recorded as
still open rather than quietly re-dated. Separately, the Philly status-split
rung, though written to specification, has never been run against the real
one-gigabyte trace in this build --- the artifact is stored with Git LFS and a
plain fetch returns a pointer file. Published targets, tolerances and structural
assertions on a small synthetic slice exist, but there is no measured Philly
number here, and the suite marks those rows as unmeasured rather than as
agreement.

\begin{table}[t]
\centering\small
\caption{Anti-goals: published results fleetsim deliberately does not
reproduce.}
\label{tab:antigoals}
\begin{tabular}{p{0.31\textwidth}p{0.42\textwidth}l}
\toprule
Published number & Why not & Disposition \\
\midrule
Philly $\sim$52.3\,\% GPU utilization \cite{philly2019} &
A hardware cycle counter, not scheduler occupancy. A discrete-event simulator
produces \emph{allocation} occupancy: a different and higher quantity. &
Out of scope \\
Philly fair-share vs.\ fragmentation-delay split \cite{philly2019} &
Needs per-job delay-cause attribution, which fleetsim does not compute. &
Still deferred \\
Helios QSSF column \cite{helios2021} &
Needs a job-name column absent from the public CSV plus a learned duration
predictor; the SJF-oracle is the reproducible upper-bound proxy. &
Out of scope \\
Alibaba 50\,\% GPU-sharing saving, median 0.042 GPU per instance
\cite{alibabapai2022} &
Needs fractional, sub-chip GPU allocation; fleetsim shares only whole chips
within a node. &
Still deferred \\
Borg absolute occupancy and JCT \cite{borg2015} &
Resources are normalized to $[0,1]$, the data is query-only, and it covers eight
independent 12{,}000-machine cells. &
Never a replay target \\
CPU contention, co-location interference &
Not modelled: an analytical speed model with no memory-bandwidth contention. &
Out of scope \\
\bottomrule
\end{tabular}
\end{table}
